% ============================================================================
% Estimation and Results
% Fragment — to be \input'd into the main document.
% Assumes: graphicx, booktabs, amsmath, subcaption (or subfig) loaded.
% ============================================================================


% ============================================================================
\section{Estimation}
\label{sec:estimation}
% ============================================================================

\subsection{Method of Simulated Moments}
\label{sec:msm}

We estimate the structural parameters of the life-cycle model using the Method
of Simulated Moments (MSM).  Let $\omega$ denote the vector of preference and
job-finding parameters to be estimated in the second stage, after
pre-estimating the stochastic processes that govern health, mortality, partner
status, wages, job separations, and care demand from auxiliary data.  The MSM
estimator chooses $\omega$ to minimize the weighted distance between moments
computed from the data and the corresponding moments computed from a sample
simulated under the model:
%
\begin{equation}
\label{eq:msm}
  \hat\omega
  \;=\;
  \arg\min_{\omega}\;
  \bigl[\hat\psi - \psi(\omega)\bigr]^{\!\top}
  \,\Sigma\,
  \bigl[\hat\psi - \psi(\omega)\bigr],
\end{equation}
%
where $\hat\psi$ is the vector of empirical moments, $\psi(\omega)$ collects
the simulated counterparts obtained by solving and simulating the life-cycle
model at parameter values~$\omega$, and $\Sigma$ is a diagonal weighting
matrix.  The diagonal elements of~$\Sigma$ are set to the inverse of the
estimated variances of the empirical moments, so that moments measured with
greater precision receive proportionally more weight in the criterion
function.\footnote{Using a diagonal weighting matrix avoids estimating the
full covariance structure of the moments, which is infeasible given the
large moment vector.  The resulting estimator is consistent, and inference
is based on a sandwich-type variance estimator that accounts for the
diagonal approximation.}

The parameter vector~$\omega$ comprises the disutility-of-work parameters
(by employment type, education, health, and caregiving status), the utility
of the caregiving arrangement (informal versus formal care, by education
and health), and the job-finding probability parameters.  Discount factor,
interest rate, wage equation coefficients, and all stochastic-process
parameters are held fixed at their pre-estimated values.

We solve the model by backward induction on a discretized state space and
simulate 10{,}000 life-cycle histories for each candidate~$\omega$.  The
criterion function~\eqref{eq:msm} is minimized using a derivative-free
optimizer.


\subsection{Moments Used in Estimation}
\label{sec:moments}

Table~\ref{tab:moments_overview} provides an overview of the 969~moments
that enter the MSM criterion.  The moments are organized into five panels.

\textit{Panel~A} targets the caregiving arrangement.  We match the age
profile of informal-care provision (any, light, and intensive) in five-year
age bins from age~40 to~70, separately by education, using data from the
German Socio-Economic Panel (GSOEP) and the Survey of Health, Ageing and
Retirement in Europe (SHARE).  We additionally match the share of pure
formal care by care intensity and education, drawn from the GSOEP Innovation
Sample, and the overall share of caregivers with high education.

\textit{Panel~B} contains labor supply shares---retirement, non-employment,
part-time, and full-time---by single year of age from 30 to~70, both pooled
and by education.  These 492~moments pin down the age profile of labor
market participation and its education gradient.

\textit{Panel~C} conditions on caregiving status.  We match labor supply
shares of caregivers (all, light, and intensive) in three-year age bins from
40 to~70, both pooled and by education, yielding 360~moments that are
central to identifying the disutility parameters associated with the double
burden of work and caregiving.

\textit{Panel~D} matches mean household wealth by five-year age bin and
education from the GSOEP, covering ages~30 to~90.  These 24~wealth moments
discipline the consumption--savings margin and the bequest parameters.

\textit{Panel~E} targets dynamic behavior.  We include year-to-year
labor-supply persistence rates (working-to-working and
not-working-to-not-working) by five-year age bin and education, as well as
caregiving persistence rates by education.  These 34~transition moments help
identify the job-finding and job-separation probabilities together with the
taste-shock scale.

\input{../tables/publication/moments_overview.tex}


% ============================================================================
\section{Results}
\label{sec:results}
% ============================================================================

This section presents the estimated structural parameters and evaluates the
model's ability to replicate the empirical moments.  We first discuss the
preference parameters---the disutility of labor, the additional disutility
incurred when combining labor with informal caregiving, the job-finding
process, and the utility derived from different caregiving
arrangements---before turning to the model fit.


% --------------------------------------------------------------------------
\subsection{Disutility of Labor}
\label{sec:results_disutil}
% --------------------------------------------------------------------------

Table~\ref{tab:disutil_params} reports the estimated disutility parameters.
Panel~A displays the parameters governing the disutility of work and
unemployment in the absence of caregiving responsibilities.

The most striking pattern is the health gradient of full-time employment.
For highly educated women in bad health, the disutility of full-time work
reaches~5.00---an order of magnitude larger than the corresponding
estimate of~0.63 for women in good health.  Among the less educated, the
gradient is qualitatively similar but less extreme, rising from~0.21 in
good health to~1.66 in bad health.  Poor health thus constitutes a
powerful deterrent to full-time employment, consistent with the empirical
regularity that health shocks precipitate labor market exit, especially
among older workers.

Part-time employment exhibits a distinct pattern.  In good health,
part-time work carries a higher disutility than full-time work for both
education groups (0.67~versus~0.21 for the less educated; 1.73~versus~0.63
for the highly educated).  This reflects the fact that healthy workers
who choose part-time forgo substantial income relative to full-time
employment; the model rationalizes their part-time choice through the
combination of household composition effects and taste heterogeneity
rather than through a low disutility of part-time hours per~se.  In bad
health, the ordering reverses: part-time becomes far less costly than
full-time, which is consistent with part-time serving as a
health-compatible mode of labor force participation.

The disutility of unemployment is moderate and somewhat higher for the
highly educated (0.59~versus~0.37), reflecting the stronger labor market
attachment and higher opportunity cost of non-employment in this group.


% --------------------------------------------------------------------------
\subsection{Disutility of Labor While Caregiving}
\label{sec:results_disutil_care}
% --------------------------------------------------------------------------

Panel~B of Table~\ref{tab:disutil_params} extends the analysis to women
who simultaneously provide informal care to a parent.  The columns
distinguish the intensity of care provision (light versus intensive) and
the caregiver's own health status.

Combining full-time work with caregiving entails substantially larger
disutility than full-time work alone.  For highly educated women in good
health, the disutility of full-time work while providing light care
is~2.55, roughly four times the non-caregiving counterpart of~0.63.
Under intensive caregiving, the figure rises further to~2.43.  The
pattern is similar for the less educated, where full-time work combined
with light care (1.37) far exceeds the non-caregiving
baseline~(0.21).  These estimates capture the double burden that
empirical studies consistently document: working full-time while also
providing hands-on care to a frail parent is a profoundly costly
combination.

Part-time work while caregiving likewise carries elevated disutility
relative to the non-caregiving case, though the increase is less
pronounced than for full-time employment.  This is consistent with
part-time schedules being more compatible with caregiving obligations.

A notable asymmetry emerges for unemployment.  Among less educated
caregivers, the disutility of unemployment while caregiving is near
zero~(0.01), far below the non-caregiving estimate of~0.37.  This
suggests that for this group, non-employment combined with caregiving
is experienced as a coherent state---essentially, the woman exits the
labor market to care for her parent---rather than as involuntary
joblessness.  For the highly educated, the disutility of unemployment
while caregiving is substantially higher~(1.27), reflecting the
greater opportunity cost of labor market withdrawal.

\input{../tables/publication/disutility_of_work_params.tex}


% --------------------------------------------------------------------------
\subsection{Job-Finding Parameters}
\label{sec:results_job_finding}
% --------------------------------------------------------------------------

Table~\ref{tab:job_offer_params} reports the estimated parameters of
the job-finding probability, specified as a logistic function of a
cubic polynomial in age (scaled by~10) and an indicator for high
education.

The negative linear coefficient on age ($-0.19$) implies that the
probability of receiving a job offer declines over the life cycle, an
effect partially offset by the positive quadratic term.  The cubic
term is small and negative, producing a job-finding probability that
decreases gradually in mid-life and more steeply after age~55.  The
positive coefficient on high education ($0.09$) implies a modest
advantage in job-finding for the more educated, consistent with the
higher employability associated with greater human capital.

\input{../tables/publication/job_offer_params.tex}


% --------------------------------------------------------------------------
\subsection{Utility of Caregiving Arrangement}
\label{sec:results_care_utility}
% --------------------------------------------------------------------------

Table~\ref{tab:care_utility_params} presents the estimated utility
associated with the three caregiving arrangements available when a
parent has a care demand: light informal care (provided by the agent),
intensive informal care, and formal care.  The reference category is
``no care''---the case in which another family member provides
informal care.  All parameters are expressed as a function of the
agent's own health and education.

The estimates reveal a pronounced preference for informal care
provision when the caregiver is in good health.  Among the highly
educated, light informal care generates a utility of~1.00, the
largest single entry in the table, reflecting a strong intrinsic
motivation to care personally for a parent with moderate needs.
Less-educated women in good health also derive substantial utility
from informal care, whether light~(0.35) or intensive~(0.85).

As the caregiver's own health deteriorates, the utility of informal
care drops sharply.  For highly educated women in bad health, the
utility of light informal care falls to~$-0.01$---effectively
zero---while intensive informal care yields only~0.02.  Among
less-educated women, the decline is less dramatic but still large:
light informal care utility falls from~0.35 to~0.06, and intensive
from~0.85 to~0.09.  The caregiver's own health status thus acts as a
powerful moderator of the willingness to provide hands-on care.

Formal care, by contrast, is consistently associated with negative
utility for healthy caregivers ($-0.34$ for the less educated, $-0.22$
for the highly educated), indicating a revealed preference against
delegating care to professional providers.  When the caregiver is in
bad health, however, formal care utility rises to a small positive
value of approximately~0.04 for both education groups.  The gap
between informal and formal care therefore narrows substantially as
the caregiver's health declines: healthy women strongly prefer
providing care themselves, but as their capacity diminishes, formal
care becomes a relatively more acceptable alternative.  This pattern
is consistent with a model in which intrinsic motivation and filial
obligation drive informal care provision under favorable health
conditions, while pragmatic considerations---notably the physical
demands of caregiving---shift preferences toward formal arrangements
when the caregiver's own health is compromised.

\input{../tables/publication/care_utility_params.tex}


% --------------------------------------------------------------------------
\subsection{Model Fit}
\label{sec:results_model_fit}
% --------------------------------------------------------------------------

We evaluate the model's external validity by comparing the simulated
moments to their empirical counterparts.  The figures below plot
simulated (solid) and observed (dashed) profiles for the key targeted
outcomes.

\subsubsection{Labor Supply: Full Sample}

Figures~\ref{fig:fit_retired_all}--\ref{fig:fit_pt_all} display the
age profiles of retirement, non-employment, full-time, and part-time
shares for the full sample (pooling caregivers and non-caregivers),
separately for low-educated (left panels) and high-educated (right
panels) women.

\begin{figure}[htbp]
  \centering
  \begin{subfigure}[t]{0.48\textwidth}
    \centering
    \includegraphics[width=\textwidth]{../figures/publication/model_fit/share_retired_all_low.pdf}
    \caption{Low education}
  \end{subfigure}
  \hfill
  \begin{subfigure}[t]{0.48\textwidth}
    \centering
    \includegraphics[width=\textwidth]{../figures/publication/model_fit/share_retired_all_high.pdf}
    \caption{High education}
  \end{subfigure}
  \caption{Model fit: retirement share (full sample).}
  \label{fig:fit_retired_all}
\end{figure}

\begin{figure}[htbp]
  \centering
  \begin{subfigure}[t]{0.48\textwidth}
    \centering
    \includegraphics[width=\textwidth]{../figures/publication/model_fit/share_non_employed_all_low.pdf}
    \caption{Low education}
  \end{subfigure}
  \hfill
  \begin{subfigure}[t]{0.48\textwidth}
    \centering
    \includegraphics[width=\textwidth]{../figures/publication/model_fit/share_non_employed_all_high.pdf}
    \caption{High education}
  \end{subfigure}
  \caption{Model fit: non-employment share (full sample).}
  \label{fig:fit_nonemp_all}
\end{figure}

\begin{figure}[htbp]
  \centering
  \begin{subfigure}[t]{0.48\textwidth}
    \centering
    \includegraphics[width=\textwidth]{../figures/publication/model_fit/share_full_time_all_low.pdf}
    \caption{Low education}
  \end{subfigure}
  \hfill
  \begin{subfigure}[t]{0.48\textwidth}
    \centering
    \includegraphics[width=\textwidth]{../figures/publication/model_fit/share_full_time_all_high.pdf}
    \caption{High education}
  \end{subfigure}
  \caption{Model fit: full-time share (full sample).}
  \label{fig:fit_ft_all}
\end{figure}

\begin{figure}[htbp]
  \centering
  \begin{subfigure}[t]{0.48\textwidth}
    \centering
    \includegraphics[width=\textwidth]{../figures/publication/model_fit/share_part_time_all_low.pdf}
    \caption{Low education}
  \end{subfigure}
  \hfill
  \begin{subfigure}[t]{0.48\textwidth}
    \centering
    \includegraphics[width=\textwidth]{../figures/publication/model_fit/share_part_time_all_high.pdf}
    \caption{High education}
  \end{subfigure}
  \caption{Model fit: part-time share (full sample).}
  \label{fig:fit_pt_all}
\end{figure}


\subsubsection{Wealth}

Figure~\ref{fig:fit_wealth} shows the mean wealth profile by five-year
age bin and education.  The model tracks the hump-shaped wealth
accumulation pattern over the life cycle and captures the level
difference between education groups.

\begin{figure}[htbp]
  \centering
  \begin{subfigure}[t]{0.48\textwidth}
    \centering
    \includegraphics[width=\textwidth]{../figures/publication/model_fit/mean_wealth_low.pdf}
    \caption{Low education}
  \end{subfigure}
  \hfill
  \begin{subfigure}[t]{0.48\textwidth}
    \centering
    \includegraphics[width=\textwidth]{../figures/publication/model_fit/mean_wealth_high.pdf}
    \caption{High education}
  \end{subfigure}
  \caption{Model fit: mean wealth by five-year age bin.}
  \label{fig:fit_wealth}
\end{figure}


\subsubsection{Labor Supply: Caregivers}

Figures~\ref{fig:fit_retired_cg}--\ref{fig:fit_pt_cg} repeat the
labor supply comparison for the subsample of informal caregivers.
These moments are critical for identification of the double-burden
parameters reported in Panel~B of Table~\ref{tab:disutil_params}.

\begin{figure}[htbp]
  \centering
  \begin{subfigure}[t]{0.48\textwidth}
    \centering
    \includegraphics[width=\textwidth]{../figures/publication/model_fit/share_retired_caregiver_low.pdf}
    \caption{Low education}
  \end{subfigure}
  \hfill
  \begin{subfigure}[t]{0.48\textwidth}
    \centering
    \includegraphics[width=\textwidth]{../figures/publication/model_fit/share_retired_caregiver_high.pdf}
    \caption{High education}
  \end{subfigure}
  \caption{Model fit: retirement share (caregivers).}
  \label{fig:fit_retired_cg}
\end{figure}

\begin{figure}[htbp]
  \centering
  \begin{subfigure}[t]{0.48\textwidth}
    \centering
    \includegraphics[width=\textwidth]{../figures/publication/model_fit/share_non_employed_caregiver_low.pdf}
    \caption{Low education}
  \end{subfigure}
  \hfill
  \begin{subfigure}[t]{0.48\textwidth}
    \centering
    \includegraphics[width=\textwidth]{../figures/publication/model_fit/share_non_employed_caregiver_high.pdf}
    \caption{High education}
  \end{subfigure}
  \caption{Model fit: non-employment share (caregivers).}
  \label{fig:fit_nonemp_cg}
\end{figure}

\begin{figure}[htbp]
  \centering
  \begin{subfigure}[t]{0.48\textwidth}
    \centering
    \includegraphics[width=\textwidth]{../figures/publication/model_fit/share_full_time_caregiver_low.pdf}
    \caption{Low education}
  \end{subfigure}
  \hfill
  \begin{subfigure}[t]{0.48\textwidth}
    \centering
    \includegraphics[width=\textwidth]{../figures/publication/model_fit/share_full_time_caregiver_high.pdf}
    \caption{High education}
  \end{subfigure}
  \caption{Model fit: full-time share (caregivers).}
  \label{fig:fit_ft_cg}
\end{figure}

\begin{figure}[htbp]
  \centering
  \begin{subfigure}[t]{0.48\textwidth}
    \centering
    \includegraphics[width=\textwidth]{../figures/publication/model_fit/share_part_time_caregiver_low.pdf}
    \caption{Low education}
  \end{subfigure}
  \hfill
  \begin{subfigure}[t]{0.48\textwidth}
    \centering
    \includegraphics[width=\textwidth]{../figures/publication/model_fit/share_part_time_caregiver_high.pdf}
    \caption{High education}
  \end{subfigure}
  \caption{Model fit: part-time share (caregivers).}
  \label{fig:fit_pt_cg}
\end{figure}
